%================================================================================
%       Safety Critical Systems Club - Data Safety Initiative Working Group
%================================================================================
%                       DDDD    SSSS  IIIII  W   W   GGGG
%                       D   D  S        I    W   W  G   
%                       D   D   SSS     I    W W W  G  GG
%                       D   D      S    I    WW WW  G   G
%                       DDDD   SSSS   IIIII  W   W   GGG
%================================================================================
%               Data Safety Guidance Document - LaTeX Source File
%================================================================================
%
% Description:
%   Data Safety Management Plan section.
%
%================================================================================
\chapter{Interactions Report} \label{bkm:interactions}

\dsiwgSectionQuote{For good ideas and true innovation, you need human interaction, conflict, argument, debate.}{Margaret Heffernan}
This section gives a suggested interactions report table of contents.
\section{Introduction}
\begin{description}
	\item[Scope and context] briefly sets the scene, describes the project;
	\item[Scenario];
	\item[Stakeholders];
	\item[Concept of operations, etc] (may be done by reference); \item[Assumptions];
	\item[References];
	\item[Abbreviations and acronyms],
\end{description}
\section{Interactions Description}
\begin{description}
\item[Description of system boundaries and interfaces]; 
\item[Producers and consumers of the data];
\item[A list of known data interactions with other entities] (may be systems, people, etc). Interactions involving data flows (e.g. messaging, networking) are of primary interest here, and control interactions (e.g based on configuration data) are also considered. It should include both external and internal data interactions to the system.
\item[Other possible interactions]
\end{description}
\section{Interactions Analysis}
The results of the failure analysis of the interactions, which may use such techniques as \gls{hazop}, data \gls{fmeca}, STAMP, FRAM, etc.
\section{Summary of Findings}
A description of the significant outcomes of the analysis and the residual risks.
\section{Conclusions}
Includes actions and recommendations to manage the risks.
